% problem_id: S001-lemma-key-nonzerodivisors
% source_id: S001 (Stacks Project)
% source_file: chow.tex
% statement_locator: chow.tex lines 1132--1144
% proof_locator: chow.tex lines 1146--1260
% env: lemma
% pairing: tight (verified)
% status: complete (statement + author proof verbatim + reason + locators)
%
\begin{lemma}
\label{lemma-key-nonzerodivisors}
Let $(A, \mathfrak m)$ be a $2$-dimensional Noetherian local ring.
Let $a, b \in A$ be nonzerodivisors.
Then we have
$$
\sum
\text{ord}_{A/\mathfrak q}(\partial_{A_{\mathfrak q}}(a, b))
=
0
$$
where the sum is over the height $1$ primes $\mathfrak q$ of $A$.
\end{lemma}

\begin{proof}
If $\mathfrak q$ is a height $1$ prime of $A$ such that $a, b$
map to a unit of $A_\mathfrak q$, then $\partial_{A_\mathfrak q}(a, b) = 1$.
Thus the sum is finite. In fact, if
$V(ab) = \{\mathfrak m, \mathfrak q_1, \ldots, \mathfrak q_r\}$
then the sum is over $i = 1, \ldots, r$.
For each $i$ we pick an extension $A_{\mathfrak q_i} \subset B_i$
as in Lemma \ref{lemma-prepare-tame-symbol} for $a, b$.
By Lemma \ref{lemma-perpare-key} with $t = ab$ and the given list of primes
we may assume we have a finite local extension $A \subset B$
with $B/A$ annihilated by a power of $ab$ and such that
for each $i$ the $B_{\mathfrak q_i} \cong B_i$.
Observe that if $\mathfrak q_{i, j}$ are the primes of $B$
lying over $\mathfrak q_i$ then we have
$$
\text{ord}_{A/\mathfrak q_i}(\partial_{A_{\mathfrak q_i}}(a, b))
=
\sum\nolimits_j
\text{ord}_{B/\mathfrak q_{i, j}}(\partial_{B_{\mathfrak q_{i, j}}}(a, b))
$$
by Lemma \ref{lemma-norm-down-tame-symbol} and
Algebra, Lemma \ref{algebra-lemma-finite-extension-dim-1}.
Thus we may replace $A$ by $B$ and
reduce to the case discussed in the next paragraph.

\medskip\noindent
Assume for each $i$ there is a nonzerodivisor
$\pi_i \in A_{\mathfrak q_i}$ and units $u_i, v_i \in A_{\mathfrak q_i}$
such that for some integers $e_i, f_i \geq 0$ we have
$$
a = u_i \pi_i^{e_i},\quad b = v_i \pi_i^{f_i}
$$
in $A_{\mathfrak q_i}$. Setting
$m_i = \text{length}_{A_{\mathfrak q_i}}(A_{\mathfrak q_i}/\pi_i)$
we have
$\partial_{A_{\mathfrak q_i}}(a, b) =
((-1)^{e_if_i}u_i^{f_i}v_i^{-e_i})^{m_i}$ by definition.
Since $a, b$ are nonzerodivisors the
$(2, 1)$-periodic complex $(A/(ab), a, b)$ has vanishing cohomology.
Denote $M_i$ the image of $A/(ab)$ in $A_{\mathfrak q_i}/(ab)$.
Then we have a map
$$
A/(ab) \longrightarrow \bigoplus M_i
$$
whose kernel and cokernel are supported in $\{\mathfrak m\}$
and hence have finite length. Thus we see that
$$
\sum e_A(M_i, a, b) = 0
$$
by Lemma \ref{lemma-compare-periodic-lengths}. Hence it suffices to show
$e_A(M_i, a, b) =
- \text{ord}_{A/\mathfrak q_i}(\partial_{A_{\mathfrak q_i}}(a, b))$.

\medskip\noindent
Let us prove this first, in case $\pi_i, u_i, v_i$ are the images
of elements $\pi_i, u_i, v_i \in A$ (using the same symbols
should not cause any confusion). In this case we get
\begin{align*}
e_A(M_i, a, b) & =
e_A(M_i, u_i\pi_i^{e_i}, v_i\pi_i^{f_i}) \\
& =
e_A(M_i, \pi_i^{e_i}, \pi_i^{f_i}) -
e_A(\pi_i^{e_i}M_i, 0, u_i) +
e_A(\pi_i^{f_i}M_i, 0, v_i) \\
& =
0 -
f_im_i\text{ord}_{A/\mathfrak q_i}(u_i) +
e_im_i\text{ord}_{A/\mathfrak q_i}(v_i) \\
& =
-m_i\text{ord}_{A/\mathfrak q_i}(u_i^{f_i}v_i^{-e_i}) =
-\text{ord}_{A/\mathfrak q_i}(\partial_{A_{\mathfrak q_i}}(a, b))
\end{align*}
The second equality holds by Lemma \ref{lemma-multiply-period-length}.
Observe that
$M_i \subset (M_i)_{\mathfrak q_i} = A_{\mathfrak q_i}/(\pi_i^{e_i + f_i})$
and
$(\pi_i^{e_i}M_i)_{\mathfrak q_i} \cong A_{\mathfrak q_i}/\pi_i^{f_i}$ and
$(\pi_i^{f_i}M_i)_{\mathfrak q_i} \cong A_{\mathfrak q_i}/\pi_i^{e_i}$.
The $0$ in the third equality comes from
Lemma \ref{lemma-powers-period-length-zero}
and the other two terms come from
Lemma \ref{lemma-length-multiplication}.
The last two equalities follow from multiplicativity of
the order function and from the definition of our tame symbol.

\medskip\noindent
In general, we may first choose $c \in A$, $c \not \in \mathfrak q_i$
such that $c\pi_i \in A$. After replacing $\pi_i$ by $c\pi_i$
and $u_i$ by $c^{-e_i}u_i$ and $v_i$ by $c^{-f_i}v_i$
we may and do assume $\pi_i$ is in $A$.
Next, choose an $c \in A$, $c \not \in \mathfrak q_i$
with $cu_i, cv_i \in A$. Then we observe that
$$
e_A(M_i, ca, cb) = e_A(M_i, a, b) - e_A(aM_i, 0, c) + e_A(bM_i, 0, c)
$$
by Lemma \ref{lemma-length-multiplication}.
On the other hand, we have
$$
\partial_{A_{\mathfrak q_i}}(ca, cb) =
c^{m_i(f_i - e_i)}\partial_{A_{\mathfrak q_i}}(a, b)
$$
in $\kappa(\mathfrak q_i)^*$ because $c$ is a unit in $A_{\mathfrak q_i}$.
The arguments in the previous paragraph show that
$e_A(M_i, ca, cb) = -
\text{ord}_{A/\mathfrak q_i}(\partial_{A_{\mathfrak q_i}}(ca, cb))$.
Thus it suffices to prove
$$
e_A(aM_i, 0, c) = \text{ord}_{A/\mathfrak q_i}(c^{m_if_i})
\quad\text{and}\quad
e_A(bM_i, 0, c) = \text{ord}_{A/\mathfrak q_i}(c^{m_ie_i})
$$
and this follows from Lemma \ref{lemma-length-multiplication}
by the description (see above)
of what happens when we localize at $\mathfrak q_i$.
\end{proof}
